\section{Operational Semantics}
\label{sec:opsem}

The operational semantics has one primitive execution relation: the
continuation-machine step $s\longrightarrow s'\,;\,\eta$.  The other judgments
below are derived or proof-facing relations built around that step.  They are
kept distinct because they play different roles in the metatheory.

\begin{center}
\small
\begin{tabular}{@{}>{\raggedright\arraybackslash}p{.18\linewidth}
                  >{\raggedright\arraybackslash}p{.33\linewidth}
                  >{\raggedright\arraybackslash}p{.39\linewidth}@{}}
\toprule
\textbf{Layer} & \textbf{Judgment} & \textbf{Role}\\
\midrule
machine step
  & $s\longrightarrow s'\,;\,\eta$
  & one continuation-machine step emitting one trace fragment\\
finite execution
  & $s\longrightarrow^{*}s'\,;\,\phi$
  & accumulated machine trace\\
trace replay
  & $(\phi,H)\Rightarrow^{*}(\Nil,H')$
  & replay of a produced trace over a heap\\
\bottomrule
\end{tabular}
\end{center}

The notation is intentionally separated.  Machine execution \emph{produces}
trace fragments; the finite-run relation concatenates them into a structured
trace; replay later \emph{consumes} an already produced trace.  A one-step
transition emits only a fragment $\eta$: administrative steps emit $\Nil$, while
heap operations emit singleton fragments such as $\Alloc(l\at r,v)$,
$\Read(l\at r,v)$, and $\Write(l\at r,v)$.  The trace of a longer execution is
not stored in the state or continuation.

\begin{formalbox}{Operational Domains}
\small
\begin{tabular}{@{}>{$}l<{$} @{\quad} l@{}}
E\in\mathit{Env}\triangleq \mathit{Variables}\rightharpoonup\mathit{Values} & (variable assignment)\\
P\in\mathit{Rgn}\triangleq \mathit{RgnVars}\rightharpoonup\mathit{RgnConst} & (regions in scope)\\
H\in\mathit{Heap}\triangleq \mathit{Addr}\rightharpoonup\mathit{Values} & (heap assignment)\\
\langle E,P,e\rangle\in\mathit{Closures}\triangleq
  \mathit{Env}\times\mathit{Rgn}\times\mathit{Expressions} & (function closures)\\
v\in\mathit{Values}\triangleq
  \mathbb{N}+\mathbb{B}+\mathit{Addr}+\mathit{Closures}+\Theta+()+ (v,v) & (dynamic values)\\[.6ex]
s ::= \langle H,E,P,e,k\rangle\mid\langle H,v,k\rangle\mid\mathsf{done}(H,v) & (machine states)\\
a ::= \Alloc(l\at r,v)\mid\Read(l\at r,v)\mid\Write(l\at r,v) & (dynamic actions)\\
\eta ::= \Nil\mid a & (one-step trace fragments)\\
\phi ::= \Nil\mid a\mid \phi\tracecat\phi\mid \phi\parallel\phi & (structured traces)
\end{tabular}
\end{formalbox}

\begin{formalbox}{Finite Machine Executions}
\footnotesize
\noindent\textbf{Finite execution}\hfill $\boxed{s\longrightarrow^{*}s'\,;\,\phi}$

Each finite execution accumulates one-step fragments by induction.
\begin{mathpar}
\inferrule*[lab=TR-Refl]{ }
  {s\longrightarrow^{*}s\,;\,\Nil}

\inferrule*[lab=TR-Step]
  {s\longrightarrow s_1\,;\,\eta \\
   s_1\longrightarrow^{*}s_2\,;\,\phi}
  {s\longrightarrow^{*}s_2\,;\,\eta\tracecat\phi}
\end{mathpar}
Here $\tracecat$ is structured trace sequencing and $\Nil$ is its identity.
Ordinary expression evaluation is the special case
\[
\mathsf{init}(H,E,P,e)
  \longrightarrow^{*}
\mathsf{done}(H',v)\,;\,\phi.
\]
A surface-effect term is just an ordinary expression of type $\mathsf{eff}$;
there is no separate summary evaluator.
\end{formalbox}

\begin{formalbox}{Trace Replay}
\footnotesize
\noindent\textbf{Replay}\hfill $\boxed{(\phi,H)\Rightarrow^{*}(\Nil,H')}$

Replay interprets an already produced structured trace as a heap transformer.
The following selected rules fix the intended direction: execution produces
traces, while replay consumes them.
\begin{mathpar}
\inferrule*[lab=R-Nil]{ }
  {(\Nil,H)\Rightarrow^{*}(\Nil,H)}

\inferrule*[lab=R-Alloc]{ }
  {(\Alloc(l\at r,v),H)\Rightarrow^{*}(\Nil,H[(r,l)\mapsto v])}

\inferrule*[lab=R-Read]{H(r,l)=v}
  {(\Read(l\at r,v),H)\Rightarrow^{*}(\Nil,H)}

\inferrule*[lab=R-Write]{H(r,l)\neq\bot}
  {(\Write(l\at r,v),H)\Rightarrow^{*}(\Nil,H[(r,l)\mapsto v])}

\inferrule*[lab=R-Seq]
  {(\phi_1,H)\Rightarrow^{*}(\Nil,H_1) \\
   (\phi_2,H_1)\Rightarrow^{*}(\Nil,H_2)}
  {(\phi_1\tracecat\phi_2,H)\Rightarrow^{*}(\Nil,H_2)}
\end{mathpar}
The replay relation for parallel traces allows the two subtraces to be consumed
by an interleaving of their replay steps; the determinism theorem later shows
that, under the computed-effect guard, terminating replays agree on the final
heap.
\end{formalbox}


The initial state for expression $e$ is
$\mathsf{init}(H,E,P,e)=\langle H,E,P,e,\mathsf{doneK}\rangle$; terminal
states $\mathsf{done}(H,v)$ have no outgoing machine step.  Continuations are
bookkeeping devices for the call-by-value evaluation order.
A frame records the expression still to be evaluated, the environment and region
map needed to resume evaluation, and the continuation to resume afterward.
Application frames first remember the unevaluated argument while the function
position is evaluated, then remember the closure while the argument is
evaluated, and finally enter either the computational body or the summary body
with the recursive name and argument bound.  Heap-operation frames remember the
pending operation and the region environment while a location or value
subexpression is evaluated.  The trace emitted while evaluating functions,
arguments, or summaries is the machine trace of that evaluation.  A computed
effect is the value returned by evaluating the corresponding summary expression.

\begin{formalbox}{Small-Step Operational Semantics}
\footnotesize
\noindent\textbf{One machine step}\hfill $\boxed{s\longrightarrow s'\,;\,\eta}$
\begin{mathpar}
\inferrule*[lab=S-Const]{ }
  {\langle H,E,P,\mathsf{const}\;n,k\rangle
   \longrightarrow
   \langle H,n,k\rangle\,;\,\Nil}

\inferrule*[lab=S-Var]{E(x)=v}
  {\langle H,E,P,x,k\rangle
   \longrightarrow
   \langle H,v,k\rangle\,;\,\Nil}

\inferrule*[lab=S-Done]{ }
  {\langle H,v,\mathsf{doneK}\rangle
   \longrightarrow
   \mathsf{done}(H,v)\,;\,\Nil}

\inferrule*[lab=S-Clos]
  {\varphi=\mu f.\lambda x.\{e_b\mid e_s\}\;\vee\;\varphi=\Lambda\rho.e}
  {\langle H,E,P,\varphi,k\rangle
   \longrightarrow
   \langle H,\langle E,P,\varphi\rangle,k\rangle\,;\,\Nil}

\inferrule*[lab=S-Capp-Fun]{ }
  {\langle H,E,P,\capp e_f\,e_a,k\rangle
   \longrightarrow
   \langle H,E,P,e_f,K^\mu_f(e_a,E,P,k)\rangle\,;\,\Nil}

\inferrule*[lab=S-Capp-Arg]
  {v_f=\langle E',P',\mu f.\lambda x.\{e_b\mid e_s\}\rangle}
  {\langle H,v_f,K^\mu_f(e_a,E,P,k)\rangle
   \longrightarrow
   \langle H,E,P,e_a,K^\mu_a(E',P',f,x,e_b,e_s,k)\rangle\,;\,\Nil}

\inferrule*[lab=S-Capp-Body]
  {E''=E'[f\mapsto v_f][x\mapsto v_a] \\
   v_f=\langle E',P',\mu f.\lambda x.\{e_b\mid e_s\}\rangle}
  {\langle H,v_a,K^\mu_a(E',P',f,x,e_b,e_s,k)\rangle
   \longrightarrow
   \langle H,E'',P',e_b,k\rangle\,;\,\Nil}

\inferrule*[lab=S-Eapp-Fun]{ }
  {\langle H,E,P,\eapp e_f\,e_a,k\rangle
   \longrightarrow
   \langle H,E,P,e_f,K^e_f(e_a,E,P,k)\rangle\,;\,\Nil}

\inferrule*[lab=S-Eapp-Arg]
  {v_f=\langle E',P',\mu f.\lambda x.\{e_b\mid e_s\}\rangle}
  {\langle H,v_f,K^e_f(e_a,E,P,k)\rangle
   \longrightarrow
   \langle H,E,P,e_a,K^e_a(E',P',f,x,e_b,e_s,k)\rangle\,;\,\Nil}

\inferrule*[lab=S-Eapp-Body]
  {E''=E'[f\mapsto v_f][x\mapsto v_a] \\
   v_f=\langle E',P',\mu f.\lambda x.\{e_b\mid e_s\}\rangle}
  {\langle H,v_a,K^e_a(E',P',f,x,e_b,e_s,k)\rangle
   \longrightarrow
   \langle H,E'',P',e_s,k\rangle\,;\,\Nil}

\inferrule*[lab=S-Alloc]
  {P(p)=r \qquad \mathsf{fresh}(H,r)=l}
  {\langle H,v,K^{ref}(p,P,k)\rangle
   \longrightarrow
   \langle H[(r,l)\mapsto v],l\at r,k\rangle\,;\,\Alloc(l\at r,v)}

\inferrule*[lab=S-Get]
  {P(p)=r \qquad H(r,l)=v}
  {\langle H,l\at r,K^{get}(p,P,k)\rangle
   \longrightarrow
   \langle H,v,k\rangle\,;\,\Read(l\at r,v)}

\inferrule*[lab=S-Set]
  {P(p)=r \qquad H(r,l)\neq\bot}
  {\langle H,v,K^{set}(p,l,P,k)\rangle
   \longrightarrow
   \langle H[(r,l)\mapsto v],(),k\rangle\,;\,\Write(l\at r,v)}
\end{mathpar}
\end{formalbox}

\paragraph{Effect-valued expression steps.}
Surface-effect terms are ordinary expressions whose values are computed effects.
Abstract summaries return region-level actions without touching the heap;
concrete summaries first evaluate their argument and then return a concrete
computed action. These rules are ordinary machine steps and therefore emit
$\Nil$.

\begin{formalbox}{Effect-Valued Expression Steps}
\footnotesize
\noindent\textbf{One machine step}\hfill $\boxed{s\longrightarrow s'\,;\,\eta}$
\begin{mathpar}
\inferrule*[lab=S-Abst-Read]{P(p)=r}
  {\langle H,E,P,\Rhat(p),k\rangle
   \longrightarrow
   \langle H,\Rs(r),k\rangle\,;\,\Nil}

\inferrule*[lab=S-Abst-Write]{P(p)=r}
  {\langle H,E,P,\What(p),k\rangle
   \longrightarrow
   \langle H,\Ws(r),k\rangle\,;\,\Nil}

\inferrule*[lab=S-Abst-Alloc]{P(p)=r}
  {\langle H,E,P,\Ahat(p),k\rangle
   \longrightarrow
   \langle H,\As(r),k\rangle\,;\,\Nil}

\inferrule*[lab=S-Conc-Read]{ }
  {\langle H,l\at r,K^{read}(k)\rangle
   \longrightarrow
   \langle H,\Read(l\at r),k\rangle\,;\,\Nil}

\inferrule*[lab=S-Conc-Write]{ }
  {\langle H,l\at r,K^{write}(k)\rangle
   \longrightarrow
   \langle H,\Write(l\at r),k\rangle\,;\,\Nil}

\inferrule*[lab=S-Concat]{ }
  {\langle H,\Theta_2,K^{concat}(\Theta_1,k)\rangle
   \longrightarrow
   \langle H,\Theta_1\cup\Theta_2,k\rangle\,;\,\Nil}

\inferrule*[lab=S-Top]{ }
  {\langle H,E,P,\Top,k\rangle
   \longrightarrow
   \langle H,\top,k\rangle\,;\,\Nil}

\inferrule*[lab=S-Empty]{ }
  {\langle H,E,P,\Pure,k\rangle
   \longrightarrow
   \langle H,\emptyset,k\rangle\,;\,\Nil}
\end{mathpar}
\end{formalbox}

Effect-summary evaluation is proved read-only from typing: if a well-typed
summary execution emits a read-only trace, the final heap is the initial heap.
A summary may \emph{denote} an allocation through an abstract allocation action
without performing an allocation while the summary itself is evaluated.
